\subsection{Problem 15. Substitution with definite integrals}
Evaluate the following definite integrals using appropriate substitutions:
\[I_{1}=\int_{0}^{1}\frac{2x}{1+x^{2}}dx, \quad I_{2}=\int_{0}^{\pi/2}\sin^{3}x \cos x dx, \quad I_{3}=\int_{1}^{e}\frac{\ln x}{x}dx.\]

\begin{enumerate}[label=(\alph*)]
    \item Evaluate $I_{1}$ and explain why a substitution of the form $u=1+x^{2}$ is natural.
    \item Evaluate $I_{2}$ and explain your choice of substitution in terms of recognizing a derivative inside the integrand.
    \item Evaluate $I_{3}$ and then differentiate your result with respect to the upper limit $e$ to check consistency with the FTC.
    \item For one of the integrals above, propose an incorrect substitution and explain why it fails or complicates the problem, illustrating the importance of a good choice.
\end{enumerate}

\subsection*{Strategy}
To evaluate these definite integrals, we look for a substitution $u = g(x)$ such that the derivative $g'(x)dx$ (or a multiple of it) appears in the integrand. The general strategy is:
\begin{enumerate}[label = (\alph*)]
    \item \textbf{Identify the inner function:} Look for a function "inside" another function (e.g., inside a denominator, a power, or a logarithm).
    \item \textbf{Check the derivative:} Verify if the derivative of that inner function is present as a multiplicative factor.
    \item \textbf{Change limits}
\end{enumerate}

\subsection*{Solution}
\textbf{(a) Evaluation of $I_1$:} \\
    We have $I_1 = \int_{0}^{1} \frac{2x}{1+x^2} dx$.
    \begin{itemize}
        \item Let $u = 1 + x^2$. Then, $du = 2x \, dx$.
        \item
        \begin{itemize}
            \item When $x = 0 \implies u = 1 + 0^2 = 1$.
            \item When $x = 1 \implies u = 1 + 1^2 = 2$.
        \end{itemize} 
        \item \textbf{We have:}
        \[
        I_1 = \int_{1}^{2} \frac{1}{u} \, du = \ln|u| \Big|_{1}^{2} = \ln(2) - \ln(1) = \ln(2).
        \]
        \item \textbf{Explanation:} The substitution $u = 1 + x^2$ is natural because the numerator $2x$ is exactly the derivative of the denominator $1 + x^2$. This transforms a rational function into a standard natural logarithm integral form $\int \frac{du}{u}$.
    \end{itemize}

\textbf{(b) Evaluation of $I_2$:} \\
    We have $I_2 = \int_{0}^{\pi/2} \sin^3 x \cos x \, dx$.
    \begin{itemize}
        \item Let $u = \sin x$. Then, $du = \cos x \, dx$.
        \item
        \begin{itemize}
            \item When $x = 0 \implies u = \sin(0) = 0$.
            \item When $x = \pi/2 \implies u = \sin(\pi/2) = 1$.
        \end{itemize}
        \item \textbf{We have:}
        \[
        I_2 = \int_{0}^{1} u^3 \, du = \frac{u^4}{4} \Big|_{0}^{1} = \frac{1^4}{4} - \frac{0^4}{4} = \frac{1}{4}.
        \]
        \item \textbf{Explanation:} We recognize that $\cos x$ is the derivative of $\sin x$. Since $\sin x$ is raised to a power, choosing $u = \sin x$ absorbs the $\cos x$ term into $du$, simplifying the integrand to a basic power rule form $u^n$.
    \end{itemize}

\textbf{(c) Evaluation of $I_3$ and FTC Check:} \\
    We have $I_3 = \int_{1}^{e} \frac{\ln x}{x} dx = \int_{1}^{e} (\ln x) \cdot \frac{1}{x} dx$.
    \begin{itemize}
        \item Let $u = \ln x$. Then, $du = \frac{1}{x} dx$.
        \item
        \begin{itemize}
            \item When $x = 1 \implies u = \ln(1) = 0$.
            \item When $x = e \implies u = \ln(e) = 1$.
        \end{itemize}
        \item \textbf{We have:}
        \[
        I_3 = \int_{0}^{1} u \, du = \frac{u^2}{2} \Big|_{0}^{1} = \frac{1}{2} - 0 = \frac{1}{2}.
        \]
        \item \textbf{FTC Check:} Consider the integral with a variable upper limit $t$: $A(t) = \int_{1}^{t} \frac{\ln x}{x} dx = \frac{(\ln t)^2}{2}$.
        Differentiating this result with respect to $t$:
        \[
        A'(t) = \frac{1}{2} \cdot 2(\ln t) \cdot (\ln t)' = (\ln t) \cdot \frac{1}{t} = \frac{\ln t}{t}.
        \]
        This matches the original integrand.
    \end{itemize}

\textbf{(d) Incorrect Substitution Example:} \\
    Consider $I_1 = \int_{0}^{1} \frac{2x}{1+x^2} dx$.
    \begin{itemize}
        \item \textbf{Proposed Substitution:} Let $u = 2x$. Then $x = u/2$ and $dx = du/2$.
        \item \textbf{Resulting Integral:}
        \[
        \int \frac{u}{1 + (u/2)^2} \frac{du}{2} = \int \frac{u}{2(1 + u^2/4)} du = \int \frac{2u}{4 + u^2} du.
        \]
        \item \textbf{Why it fails/complicates:} While not mathematically "wrong" (it is still solvable), this substitution does not simplify the structure of the problem. 
    \end{itemize}
\pagebreak